% GENERATED FILE. Edit canonical lessons, then run npm run generate:books.
% canonical-source-hash: fnv1a64:c4259dd2385e5b51
% canonical-lessons: KA-C141-hagura, KA-C141-mrdu, KA-C141-gatti, KA-C141-kolaku, KA-C141-vega

\chapter{Light, Soft, Hard, Dirty, Speed}
\label{ch:ka-light-soft-hard-dirty-speed}
\hlchaptermodality{141}

\begin{chapteropening}
\textbf{By the end of this chapter:} \emph{I can say light, soft, hard, dirty and speed.}

Complete the last lesson of chapter 141: i can say light, soft, hard, dirty and speed.
\end{chapteropening}

\section[hagura]{\kn{ಹಗುರ} (hagura) — light (in weight)}

\label{lesson:KA-C141-hagura}



\begin{quote}
Before the new one: say the Kannada for short, then the Kannada for heavy.
\end{quote}

\subsection*{You'll want to know: \kn{ಹಗುರ}}
\textbf{\kn{ಹಗುರ}} — \emph{hagura} — \textquotedblleft{}light (in weight)\textquotedblright{}.

\begin{grammarlens}[title={What you've built}]
One more word for this chapter.
\end{grammarlens}

\subsection*{Your turn}
\begin{itemize}
\raggedright
  \item \emph{Say it:} \emph{hagura}
  \item \emph{Say it:} \emph{hagura}, once more
  \item \emph{Say it:} say \emph{hagura}
  \item \emph{Recall:} say the Kannada for short, then the Kannada for heavy, then say \emph{hagura} again
\end{itemize}

\begin{tcolorbox}[breakable,colback=teal!4,colframe=teal!35!black,title={Before you move on}]
What does \kn{ಹಗುರ} mean? (\textquotedblleft{}Light (in weight)\textquotedblright{}.) Say it once more.
\end{tcolorbox}
\section[mṛdu]{\kn{ಮೃದು} (mṛdu) — soft}

\label{lesson:KA-C141-mrdu}



\begin{quote}
Before the new one: say the Kannada for heavy, then the Kannada for light.
\end{quote}

\subsection*{You'll want to know: \kn{ಮೃದು}}
\textbf{\kn{ಮೃದು}} — \emph{mṛdu} — \textquotedblleft{}soft\textquotedblright{}.

\begin{grammarlens}[title={What you've built}]
One more word for this chapter.
\end{grammarlens}

\subsection*{Your turn}
\begin{itemize}
\raggedright
  \item \emph{Say it:} \emph{mṛdu}
  \item \emph{Say it:} \emph{mṛdu}, once more
  \item \emph{Say it:} say \emph{mṛdu}
  \item \emph{Recall:} say the Kannada for heavy, then the Kannada for light, then say \emph{mṛdu} again
\end{itemize}

\begin{tcolorbox}[breakable,colback=teal!4,colframe=teal!35!black,title={Before you move on}]
What does \kn{ಮೃದು} mean? (\textquotedblleft{}Soft\textquotedblright{}.) Say it once more.
\end{tcolorbox}
\section[gaṭṭi]{\kn{ಗಟ್ಟಿ} (gaṭṭi) — hard, strong}

\label{lesson:KA-C141-gatti}



\begin{quote}
Before the new one: say the Kannada for light, then the Kannada for soft.
\end{quote}

\subsection*{You'll want to know: \kn{ಗಟ್ಟಿ}}
\textbf{\kn{ಗಟ್ಟಿ}} — \emph{gaṭṭi} — \textquotedblleft{}hard, strong\textquotedblright{}.

\begin{grammarlens}[title={What you've built}]
One more word for this chapter.
\end{grammarlens}

\subsection*{Your turn}
\begin{itemize}
\raggedright
  \item \emph{Say it:} \emph{gaṭṭi}
  \item \emph{Say it:} \emph{gaṭṭi}, once more
  \item \emph{Say it:} say \emph{gaṭṭi}
  \item \emph{Recall:} say the Kannada for light, then the Kannada for soft, then say \emph{gaṭṭi} again
\end{itemize}

\begin{tcolorbox}[breakable,colback=teal!4,colframe=teal!35!black,title={Before you move on}]
What does \kn{ಗಟ್ಟಿ} mean? (\textquotedblleft{}Hard, strong\textquotedblright{}.) Say it once more.
\end{tcolorbox}
\section[koḷaku]{\kn{ಕೊಳಕು} (koḷaku) — dirty}

\label{lesson:KA-C141-kolaku}



\begin{quote}
Before the new one: say the Kannada for soft, then the Kannada for hard.
\end{quote}

\subsection*{You'll want to know: \kn{ಕೊಳಕು}}
\textbf{\kn{ಕೊಳಕು}} — \emph{koḷaku} — \textquotedblleft{}dirty\textquotedblright{}.

\begin{grammarlens}[title={What you've built}]
One more word for this chapter.
\end{grammarlens}

\subsection*{Your turn}
\begin{itemize}
\raggedright
  \item \emph{Say it:} \emph{koḷaku}
  \item \emph{Say it:} \emph{koḷaku}, once more
  \item \emph{Say it:} say \emph{koḷaku}
  \item \emph{Recall:} say the Kannada for soft, then the Kannada for hard, then say \emph{koḷaku} again
\end{itemize}

\begin{tcolorbox}[breakable,colback=teal!4,colframe=teal!35!black,title={Before you move on}]
What does \kn{ಕೊಳಕು} mean? (\textquotedblleft{}Dirty\textquotedblright{}.) Say it once more.
\end{tcolorbox}
\section[vēga]{\kn{ವೇಗ} (vēga) — speed; fast}

\label{lesson:KA-C141-vega}



\begin{quote}
Before the new one: say the Kannada for hard, then the Kannada for dirty.
\end{quote}

\subsection*{You'll want to know: \kn{ವೇಗ}}
\textbf{\kn{ವೇಗ}} — \emph{vēga} — \textquotedblleft{}speed; fast\textquotedblright{}.

\begin{grammarlens}[title={What you've built}]
One more word for this chapter.
\end{grammarlens}

\subsection*{Your turn}
\begin{itemize}
\raggedright
  \item \emph{Say it:} \emph{vēga}
  \item \emph{Say it:} \emph{vēga}, once more
  \item \emph{Say it:} say \emph{vēga}
  \item \emph{Recall:} say the Kannada for hard, then the Kannada for dirty, then say \emph{vēga} again
\end{itemize}

\begin{tcolorbox}[breakable,colback=teal!4,colframe=teal!35!black,title={Before you move on}]
What does \kn{ವೇಗ} mean? (\textquotedblleft{}Speed; fast\textquotedblright{}.) Say it once more.
\end{tcolorbox}
